% Einheit 1 von 5 – Scharbegriff und Parameterwert (bank/funktionsscharen-und-ortskurven/e1.jsonl, zusammenbau v0.1)
%% TODO Zeitmarke Einheit 1: Marken-Zeile nicht lesbar
%% TODO Zweigzeile Teil 1: was hier gelernt wird (Satz in Schülersprache aus dem Einheitstitel) – Einheit 1
\einheitenkopf[e1]{Einheit 1 von 5 · Scharbegriff und Parameterwert}
\zweigzeile{Abitur GK}
\setcounter{aufgabe}{9}

%% TODO Titel: Ich-kann-Satz fehlt in der Bank (Platzhalter: Kettenname) – Nr. 10
%% TODO Anweisung: ein Satz über der ersten Teilaufgabe fehlt in der Bank – Nr. 10
\begin{aufgabe}{Scharbegriff und Parameterwert}
%% TODO \rechenplatz{n}: Schrittzahl des Grundfalls fehlt in der Bank (nur bei fester Schreibform, 2.3 b)
\begin{teile}
\teil Gegeben ist die Schar $f_a(x) = a x^2 + 3$ mit dem Parameter $a$. Gesucht ist der Wert von $a$, für den der Graph durch $P(2 | 7)$ geht. Wonach wird aufgelöst? \\ \kreuz{nach $x$} \kreuz{nach $a$}
\teil Gegeben ist die Schar $f_a(x) = a x^2 - 2x$ mit $a \in \mathbb{R}$. Schreibe $f_3(x)$ auf und berechne $f_3(2)$. $f_3(x) =$ \_\_, $f_3(2) =$ \_\_
\teil Gegeben ist die Schar $f_a(x) = a x^4 + a x$ mit $a > 0$. Für welchen Wert von $a$ liegt $P(1 | 10)$ auf dem Graphen von $f_a$? \hfill (Abitur 2024 LK) \\ $a =$ \_\_
\teil Gegeben ist die Schar $f_a(x) = 2x^2 \cdot e^{a x}$ mit $a \neq 0$. Der Graph von $f_a$ geht durch $P(-3 \,|\, 18e^{-6})$. Bestimme $a$. \hfill (Abitur 2025 LK) \\ $a =$ \_\_
\teil Eine Rinne hat im Querschnitt das Profil $p_a(x) = a x^2$ mit $a > 0$ ($x$ und $p_a(x)$ in dm). Sie ist oben 4 dm breit und 2 dm tief. Bestimme $a$. $a =$ \_\_
\teil Die Abbildung zeigt die Graphen I, II und III der Schar $f_c(x) = x^2 - 2x + c$ für $c = -2$, $c = 0$ und $c = 2$. Ordne zu.

\begin{ksys}[xmin=-2,xmax=4,ymin=-4,ymax=6,ablesen] \funktionab{\x^2-2*\x}{I}{-2}{4} \funktionab{\x^2-2*\x+2}{II}{-1.2}{3.2} \funktionab{\x^2-2*\x-2}{III}{-2}{4} \end{ksys}
I: \leerfeld , II: \leerfeld , III: \leerfeld
\teil Gegeben ist die Schar $f_a(x) = (a - 3)x^2 + 2x + a$ mit $a \in \mathbb{R}$. Für welches $a$ ist $f_a$ linear? Gib Steigung und y-Achsenabschnitt dieses Sonderfalls an. $a =$ \_\_, Steigung \_\_, y-Achsenabschnitt \_\_
\teil Die Abbildung zeigt den Graphen von $p_{0{,}5}$ der Schar $p_k(x) = k x^2 - 1$ mit $k > 0$. Skizziere den Graphen von $p_2$ in die Abbildung. \hfill (Abitur 2022 LK)

\begin{ksys}[xmin=-3,xmax=3,ymin=-2,ymax=8,ablesen] \funktionab{0.5*\x^2-1}{p_{0{,}5}}{-3}{3} \end{ksys}
\teil Die Höhe einer Sonnenblume in cm wird durch $H(t) = \frac{p}{q + e^{r t}}$ mit $p, q > 0$ und $r < 0$ modelliert ($t$ in Tagen). Für eine Pflanze gilt: (1) $\frac{p}{q + 1} = 10$; (2) für $t \to \infty$ strebt $H(t)$ gegen $200$; (3) $\frac{p}{q + e^{30r}} = 100$. Deute jede Gleichung im Sachzusammenhang. \hfill (Abitur 2024 LK)
\end{teile}
\end{aufgabe}

%% TODO Titel: Ich-kann-Satz fehlt in der Bank (Platzhalter: Kettenname) – Nr. 11
%% TODO Anweisung: ein Satz über der ersten Teilaufgabe fehlt in der Bank – Nr. 11
\begin{aufgabe}{Scharbegriff und Parameterwert – Fehler finden, Begründen, Anwendung}
\begin{teile}
\teil Gesucht ist $k$ so, dass $Q(2 | 9)$ auf dem Graphen von $g_k(x) = k x^2 - 3$ liegt. Tim rechnet: \rechnung{2x^2 - 3 &= 9 \\ x^2 &= 6 \\ x &= \pm\sqrt{6}} Finde den Fehler und rechne richtig.
\teil Begründe, warum zu $a = 2$ genau ein Graph der Schar $f_a(x) = a x^2 + 1$ gehört.
\teil Ein Brückenbogen folgt der Kurve $b_a(x) = -a x^2 + 9$ mit $a > 0$ ($x$ und $b_a(x)$ in m) über dem Wasserspiegel $y = 0$. Am Wasser soll der Bogen 12 m breit sein. Bestimme $a$ und die Höhe des Bogens 3 m neben der Mitte. $a =$ \_\_, Höhe: \_\_ m
\end{teile}
\end{aufgabe}
